\clearpage
\subsection{Sudeste asiático: ASEAN}

La \acrfull{asean} se estableció en 1967, enmarcada por la guerra de Vietnam, por 5 países: Malasia, Indonesia, Filipinas, Singapur y Tailandia. Su objetivo principal era promover la paz y la estabilidad en la región, así como fomentar la cooperación económica, social, cultural, técnica y educativa entre los países miembro. Posteriormente, se unieron a la asociación otros 6 países: Brunei Darussalam en 1984, Vietnam en 1995, Laos y Myanmar en 1997, Camboya en 1999 y, recientemente, en 2025, Timor Oriental \parencite{ASEAN}. Durante los primeros años de existencia, la cooperación económica en la \gls{asean} fue promovida a través de un \gls{pta} que tenía como uno de sus principales objetivos reducir los aranceles sobre ciertos productos industriales. Sin embargo, este objetivo no logró materializarse debido, principalmente, a que los países integrantes de la \gls{asean} siguieron una estrategia de industrialización por sustitución de importaciones durante esos años. La integración económica de la \gls{asean} comenzó formalmente con la creación del \gls{afta} en 1992. Desde entonces, el proceso de integración regional ha pasado por tres fases principales: una etapa inicial caracterizada por la consolidación del \gls{afta} de 1992 a 2002; después, una fase marcada por la conformación de la \gls{aec} entre 2003 y 2015; y una tercera etapa centrada en la implementación del Plan Maestro de la \gls{aec} 2025 a partir de 2016 \parencite[pp.~25--26]{Ishikawa2021}.

El \gls{afta} fue creado como respuesta al surgimiento de otros bloques regionales, como el \gls{tlcan} y la expansión de la \gls{ue}. Sus objetivos principales fueron la creación de un mercado único, la atracción de inversiones extranjeras directas en la región y la expansión del comercio y las inversiones intrarregionales. El principal mecanismo para lograr estos objetivos fue el \gls{cept}, el cual permitió una reducción gradual de los aranceles dentro de la \gls{asean} hasta alcanzar una tasa arancelaria entre el 0 y el 5\%. Para acceder a estas preferencias, los bienes debían cumplir reglas de origen, basadas en un contenido regional acumulado de al menos 40\% \parencite{ASEAN_CEPT}. Sin embargo, a diferencia de la \gls{ue}, el \gls{afta} no implicó un arancel externo común a las mercancías importadas de fuera de la asociación \parencite{ASEAN_AFTA}. El \gls{afta} logró que los aranceles en la \gls{asean}6 (Brunei, Indonesia, Malasia, Filipinas, Singapur y Tailandia) se redujeran a menos del 5\% en 2003, cinco años antes de lo planificado, y se eliminaran en prácticamente todas las líneas arancelarias en 2010. \textcite{Ishikawa2021} sostiene que el éxito del \gls{afta} se debió en gran medida a que se trató de un proceso gradual y flexible de liberalización, que consideró los diferentes niveles de desarrollo económico e industrial entre los países miembro. Este proceso de liberalización gradual y flexible fue favorecido por el diseño institucional del \gls{afta}, que contempló calendarios diferenciados, listas de exclusión y un trato especial para los países con menores niveles de desarrollo. El esquema \gls{cept}-\gls{afta} fue notificado ante el \gls{gatt} bajo la Cláusula de Habilitación, la cual permite acuerdos comerciales preferenciales entre países en desarrollo con requisitos menos estrictos que los establecidos en el artículo XXIV para las zonas de libre comercio y las uniones aduaneras.

La integración económica de la \gls{asean} se amplió durante la década de 1990, más allá de la liberalización comercial. En 1995 se estableció el \gls{afas}, orientado a reducir las restricciones al comercio regional de servicios \parencite{ASEAN_services}. Al año siguiente, se creó el esquema \gls{aico}, que promovió las actividades industriales conjuntas y la formación de redes regionales de producción; y en 1998 se constituyó el Área de Inversión de la \gls{asean}, y después fue sustituida por el \gls{acia} en 2009, con el propósito de liberalizar, facilitar y proteger las inversiones intrarregionales \parencite[p.~16]{ASEAN_AIR2012}. En conjunto, estos instrumentos fortalecieron el comercio, la inversión extranjera y la consolidación de cadenas regionales de valor.

La \acrfull{aec} fue planteada por primera vez en 2003 durante la IX Cumbre de la asociación, con el objetivo de transformar a la \gls{asean} en una región con libre circulación de bienes, servicios, inversiones, mano de obra calificada y una mayor libertad para los flujos de capital. Como parte de esta estrategia, se firmó el Acuerdo Marco de la \gls{asean} para la Integración de Sectores Prioritarios, que incluyó los productos agroindustriales, las industrias automotriz, electrónica y textil, la pesca, los productos derivados del caucho y la madera, el transporte aéreo, las tecnologías de la información y la comunicación, los servicios de salud, el turismo y la logística. En la XII Cumbre de la \gls{asean} celebrada en 2007, se aceleró el compromiso de establecer la \gls{aec} para 2015. Estos acuerdos también contemplaban el mejoramiento de la infraestructura de transporte y comunicaciones intrarregional, así como el establecimiento de mecanismos de interconexión energética \parencite{Ishikawa2021, ASEAN_Vision_2020}.

La \gls{aec} fue formalmente establecida en 2015, conforme al calendario previsto, aunque no todas las medidas acordadas en el Plan Maestro de la \gls{aec} 2015 fueron implementadas. Un elemento clave que ayudó al cumplimiento de este compromiso fue la sustitución del esquema \gls{cept} por el \gls{atiga} en 2010. El logro más importante del \gls{aec} fue la eliminación de aranceles, alcanzando una tasa de liberalización comercial\footnote{Porcentaje de líneas arancelarias sujetas a eliminación} del 98.6\% \parencite[p.~29]{Ishikawa2021}. El ATIGA también abarcó la facilitación de trámites aduaneros y el reconocimiento mutuo de normas \parencite[p.~16]{ASEAN_AIR2012}. En contraste, existen otros elementos que han avanzado con mayor lentitud, como la reducción de barreras no arancelarias y la liberalización del comercio de servicios y de las inversiones \parencite[p.~29]{Ishikawa2021}.

Aunque los marcos de integración económica de la \gls{asean} no incluyen un arancel común al exterior, la Asociación ha establecido de manera conjunta acuerdos de libre comercio y de asociación económica con otros países asiáticos bajo el esquema ASEAN$+1$. En particular, se establecieron acuerdos de libre comercio con la República Popular China, Corea del Sur, Japón, Australia y Nueva Zelanda, India y Hong Kong \parencite{ASEAN_FTA}. A su vez, en 2019 concluyeron las negociaciones sobre la conformación de la \gls{rcep}, un acuerdo de libre comercio multilateral entre los Estados miembros de la \gls{asean} y China, Japón, Corea del Sur, Australia y Nueva Zelanda. El objetivo del \gls{rcep} es armonizar los acuerdos de libre comercio existentes de la ASEAN$+1$ en un acuerdo unificado, con un conjunto único y coherente de normas comerciales para la región \parencite[pp.~92--94]{ASEAN_Digital2025}.
% ASEAN-CHINA FTA (ACFTA)effect 2004
% ASEAN-Korea FTA (AKFTA) effect 2007
% ASEAN-Japon FTA (AJFTA) effect 2008
% ASEAN-Australia-New Zealand FTA (AANZAFTA) effect 2010
% ASEAN-INDIA (AIFTA) effect 2010
% ASEAN-Hong-Ling (AHFKTA), 2019

%Después del establecimiento de la \gls{aec}, la integración económica regional se profundizó mediante nuevos mecanismos de facilitación comercial, apertura de servicios y vinculación externa. La \gls{asean} Single Window fue desarrollada para permitir el intercambio electrónico de documentos comerciales y aduaneros, agilizando el despacho de mercancías entre los países miembros; el \gls{atisa} amplió los compromisos regionales para reducir las barreras regulatorias al comercio de servicios; y el \gls{rcep} extendió la integración hacia los principales socios económicos de la región mediante disposiciones sobre bienes, servicios, telecomunicaciones y comercio electrónico. En conjunto, estos instrumentos contribuyeron a reducir los costos de las operaciones transfronterizas y a fortalecer la participación de la \gls{asean} en las cadenas regionales de valor \parencite[pp.~94--95]{ASEAN_Digital2025}

% 
% En 2018, la tasa de liberalización alcanzó el 99.3\% para el ASEAN6 y el 98.6\% para el ASEAN en su conjunto. \cite{Ishikawa2021} 
% Al año siguiente, durante la X Cumbre se firmó el Acuerdo marco de la ASEAN para la Integración de Sectores Prioritarios, que comprendía el periodo de 2004 a 2010, designando como sectores prioritarios los productos agropecuarios, pesca, industriales (automotriz, productos de caucho, textiles y prendas de vestir, de madera), transporte aéreo, telecomunicaciones, sanidad y turismo; previendo su plena integración para 2010 \parencite[p.~28]{Ishikawa2021}.

% La CEA aspira a consolidar a la ASEAN como un mercado único y una base de producción, una región económica altamente competitiva, una región de desarrollo económico equitativo y una región plenamente integrada en la economía global (nota 17). Estos cuatro objetivos constituyen los cuatro pilares del Plan de la CEA.

Por su parte, el análisis de las comunidades en la red insumo--producto mundial ofrece una oportunidad para evaluar el nivel de integración económica de la \gls{asean} a partir de sus vínculos intersectoriales. En primer lugar, los resultados indican un fuerte nivel de estabilidad en la integración económica de la región, en el sentido de que las economías de la mayoría de países miembros pertenecen al mismo bloque económico durante el periodo estudiado (véase la figura \ref{ComunidadesPorPais_ASEAN}). En segundo lugar, el análisis de comunidades permitió identificar a los países con menor cohesión en la economía regional de la \gls{asean}, a partir del número de años en los que sus economías conforman una comunidad individual: Indonesia (14 años), Myanmar (13 años), Brunei (6 años) y Filipinas (4 años). En contraste, también se observaron vínculos de integración con economías externas al \gls{asean}, particularmente Taiwán y Hong Kong. El caso de Taiwán resulta relevante, ya que durante todo el periodo estudiado formó parte del bloque económico asociado a la mayoría de los países de la \gls{asean}. Finalmente, los resultados no reflejan una asociación directa con el establecimiento de acuerdos de libre comercio formales. Por ejemplo, aunque la \gls{asean} ha formalizado áreas de libre comercio con países como China o India, esto no se ve reflejado directamente en la integración de sus economías en la misma comunidad. En el caso de Hong Kong, su integración al bloque económico asociado a la \gls{asean} ocurre años antes de la firma del acuerdo en 2019. Solo en el caso de Corea del Sur, las economías de Vietnam y Camboya presentan una integración estable a su comunidad a partir de 2016. Por el contrario, aunque la \gls{asean} no cuenta con un tratado de libre comercio con Taiwán, su economía se encuentra estrechamente vinculada a la de los países de la \gls{asean} a través de sus flujos intersectoriales.

\begin{figure}[H]
	\centering
	\includegraphics[width=\textwidth]{images/ComunidadesPorPais_ASEAN.png}
	\caption{Integración de las economías de la ASEAN en bloques productivos.$^{*}$}
	\label{ComunidadesPorPais_ASEAN}
	
	\caption*{\footnotesize $^{*}$ El color en la gráfica depende del país al que pertenece el vértice con mayor fuerza de salida intrarregional en ese año. Eso puede causar un efecto visual de contraste en una misma comunidad. Como muestra de ello, en 2005 y 2006 todos los países de la \gls{asean} más Taiwán formaron una sola comunidad, solo que en 2005 el nodo con mayor fuerza de salida era de Taiwán y en 2006 de Indonesia, lo cual explica el cambio abrupto de color. Por lo tanto, la gráfica no debe evaluarse únicamente por el color, sino por los países que pertenecen a un mismo color.}
\end{figure}
%\begin{figure}[H]
%	\centering
%	\includegraphics[width=\textwidth]{images/ComunidadesPorPais_ASEAN.png}
%	\caption{Integración de las economías de la ASEAN en bloques comerciales$^{*}$.}
%	\label{ComunidadesPorPais_ASEAN}

%	\caption*{\footnotesize $^{*}$ Las filas representan países y las columnas años. El color de cada celda corresponde al color de la comunidad a la cual pertenece la mayoría de los sectores económicos de un país en el año correspondiente. Cada comunidad se identifica con el país del vértice con la mayor fuerza de salida.}
%\end{figure}


La dinámica de las comunidades identificadas permite establecer cuatro periodos en el proceso de integración económica del sudeste asiático:
\begin{itemize}
	\item Formación del núcleo productivo regional (1995--2004).
	%un proceso de fortalecimiento del núcleo productivo regional, con incorporaciones y separaciones temporales. 
	Durante este periodo, el bloque económico asociado al \gls{asean} se compone de las economías de Vietnam, Camboya, Malasia, Filipinas, Singapur, Tailandia, Laos y Taiwán. Por su parte, Indonesia, Brunei y Myanmar forman comunidades individuales en algunos lapsos de años. Este periodo coincide con la consolidación del \gls{afta} y la reducción gradual de aranceles en la región. Además, en estos años hay muchos sectores individuales de los países de la \gls{asean} que forman parte de otras comunidades externas a la región, principalmente en las comunidades de Hong Kong y Japón (véase la figura \ref{DiagramaSankey_ASEAN} y el apéndice \ref{Apendice3}).\footnote{Dentro de este conjunto, destacan la extracción de minerales metálicos (B07) de Filipinas que perteneció a la comunidad de Japón de 1995 a 1998; la extracción de petróleo crudo y gas natural (B06) de Brunei que estuvo integrada a la comunidad de Japón en los años 2001 y 2003; y las actividades postales y de mensajería (H53) de varios países de la \gls{asean} que durante todo el periodo estuvieron en la comunidad encabezada por Hong Kong.}
	
	\item Máxima cohesión regional (2005--2008). En estos años, los 10 países de la \gls{asean} forman parte de la misma comunidad, junto con Taiwán; y en 2007 y 2008, también junto a Hong Kong. Este periodo se encuentra enmarcado por los trabajos de la \gls{asean} de establecer su \acrfull{aec}. 
	
	\item Reconfiguración e inestabilidad (2009--2015). El bloque económico asociado al \gls{asean} se reduce, con la salida de Myanmar, Indonesia y Filipinas de la comunidad. Esto podría reflejar las consecuencias que tuvo la crisis internacional de 2008 en el proceso de integración regional.
	 
	\item Fragmentación y formación de bloques subregionales (2016--2022). Se observa la integración estable de las economías de Vietnam y Camboya a la comunidad de Corea del Sur. Durante este periodo, nuevamente hay sectores individuales de los países de la \gls{asean} que se integran a otras comunidades, como las de China, Estados Unidos, Reino Unido y Japón (véase la figura \ref{DiagramaSankey_ASEAN} y el apéndice \ref{Apendice3}).\footnote{Por ejemplo, la extracción de minerales metálicos (B07) de Filipinas estuvo alineada con la economía de China a partir de 2017; y la industria de productos de madera (C16) y de caucho y plástico (C22) de Camboya se integró al bloque productivo estadounidense desde 2020.}
\end{itemize}

\begin{figure}[H]
	\centering
	\includegraphics[width=1\textwidth]{sankey_diagrams/asean_all.png}
	\caption{Evolución de la pertenencia a comunidades de los sectores económicos de los países de la \gls{asean}.}
	\label{DiagramaSankey_ASEAN}
\end{figure}

La figura \ref{Mapa_Comunidades_ASEAN} muestra el mapa de las comunidades en el sudoeste asiático en tres momentos: 1995, 2008 y 2022. Destaca la intensidad de los flujos comerciales entre Taiwán, Tailandia, Malasia y Singapur, los cuales abarcan a la mayoría de los sectores económicos de estos países. Además, dichos flujos presentan una distribución relativamente equilibrada entre los distintos socios comerciales. Esto refleja la profunda interdependencia económica que se ha formado entre estas economías.

\begin{figure}[H] 
	\centering
	\begin{subfigure}{.9\textwidth}
		\centering
		\includegraphics[width=\textwidth]{images/maps/ASEAN_1995_1.png}
		\caption{1995}
	\end{subfigure}
	\begin{subfigure}{.9\textwidth}
		\centering
		\includegraphics[width=\textwidth]{images/maps/ASEAN_2008_1.png}
		\caption{2008}
	\end{subfigure}
	\begin{subfigure}{.9\textwidth}
		\centering
		\includegraphics[width=\textwidth]{images/maps/ASEAN_2022_1.png}
		\caption{2022}
	\end{subfigure}
	\caption{Mapa de las comunidades identificadas en el sudeste asiático.$^{*}$}
	\label{Mapa_Comunidades_ASEAN}
	
	\caption*{\footnotesize $^{*}$ Los sectores económicos se disponen circularmente alrededor del centro del país al que pertenecen y se ordenan en el sentido de las manecillas del reloj. El color de cada sector indica la comunidad a la que pertenece. Las aristas representan únicamente los flujos comerciales cuya magnitud supera el percentil 50 de la distribución regional, y su grosor es proporcional a la intensidad de cada flujo.}
\end{figure}

%Extraccion de mienrales metalicos Filipinas - Japon 1995, 1996, 97, 98


%	 En segundo lugar, el análisis de comunidades muestra que existe una coincidencia entre las fases de integración de la ASEAN que propone \textcite{Ishikawa2021} y la evolución de las comunidades en la región. De 1995 a 2003 se observa que los países de la ASEAN formaron parte de una comu



%\textcolor{blue}{Observaciones: 1) Hay cierta estabilidad en el sentido que la mayoría de los paises de la asean siempre pertenecen a la misma comunidad, aunque la comunidad esté encabezada por otros países. 2) De 1995 a 2002 esta comunidad estuvo encabezada por Taiwán, posteriormentemaoritariamente por indonesia, y luego por Singapur. 3). Indonesi, Bruneii y miarmar son paises que por temporatadas se salen para formar su propia comunidad. Filipinas de 2011 a 2014.}


%Área de Libre Comercio de la ASEAN (AFTA, por sus siglas en ingles):
%Revisar \url{https://fta.miti.gov.my/index.php/pages/view/asean-afta}


%\begin{figure}[H]
%	\centering
%	\begin{subfigure}{.6\textwidth}
	%		\centering
	%		\includegraphics[width=\textwidth]{images/communities_membership/IDN.png}
	%		\caption{Comunidad encabezada por Indonesia}
	%	\end{subfigure}
%	\par\medskip
%	\begin{subfigure}{.6\textwidth}
	%		\centering
	%		\includegraphics[width=\textwidth]{images/communities_membership/TWN.png}
	%		\caption{Comunidad encabezada por Taiwán}
	%	\end{subfigure}
%	\par\medskip
%	\begin{subfigure}{.6\textwidth}
	%		\centering
	%		\includegraphics[width=\textwidth]{images/communities_membership/HKG.png}
	%		\caption{Comunidad encabezada por Hong Kong}
	%	\end{subfigure}
%	\caption{Cantidad de vértices por comunidades del sudeste asiático}
%\end{figure}

